% Elaborado por Fabian Gomez
% Version 1.0

\section{Stakeholders, Users, and Actors}

Esta sección identifica a las partes interesadas (stakeholders), clases de usuarios y actores externos relevantes para el Rover Olympus, con el objetivo de (i) asegurar que los requisitos RF/RNF/CON se derivan de necesidades reales, (ii) delimitar responsabilidades operacionales, y (iii) mantener coherencia entre ConOps, Casos de Uso y arquitectura MBSE/SysML.
En el contexto de este proyecto, el ``usuario'' principal del sistema es un operador/planificador en Tierra (Rover Planner / Ground Operator) que configura metas de exploración, supervisa telemetría y evalúa resultados en campañas análogas, mientras que el sistema ejecuta UC‑01 de forma autónoma.

\subsection{Stakeholder Profiles}

\textbf{Stakeholders primarios (directamente responsables del éxito del proyecto):}
\begin{enumerate}
    \item \textbf{SETEC Lab / Escuela de Ingeniería Electrónica (TEC):} Define el contexto académico, objetivos de validación, criterios de aceptación y restricciones de recursos.
    \item \textbf{Investigador Principal / Responsable de Ingeniería de Sistemas:} Asegura la coherencia entre ConOps, requisitos, arquitectura y V\&V.
    \item \textbf{Equipo del proyecto ELANav} (\emph{Biblioteca embebida para la navegación de sistemas espaciales de misión crítica}, TEC, 2025--2027; coordinador: Johan Carvajal Godínez): Usuario final de la validación que entrega Olympus; define los requisitos de la biblioteca embebida que el rover debe poder ejecutar y exponer mediante telemetría/comandos.
\end{enumerate}

\textbf{Stakeholders secundarios (impacto indirecto):}
\begin{enumerate}[resume]
    \item \textbf{Equipo de Integración y Pruebas (I\&T):} Responsable de la integración física, preparación de campañas y recolección de evidencia en terreno.
    \item \textbf{Evaluadores académicos / Comité de revisión:} Requieren trazabilidad, consistencia documental y evidencia reproducible de validación.
\end{enumerate}

\subsection{User Classes and Characteristics}

\textbf{Clase de usuario U1 – Rover Planner / Operador en Tierra (Ground Operator):}
\begin{itemize}
    \item \textbf{Rol:} Define objetivos de alto nivel, selecciona escenarios, configura parámetros y supervisa telemetría.
    \item \textbf{Características:} Usuario técnico familiarizado con navegación robótica; toma decisiones estratégicas sin teleoperación continua.
    \item \textbf{Necesidades:} Visibilidad del estado del sistema, evidencia de desempeño y capacidad de recuperación de datos.
\end{itemize}

\textbf{Clase de usuario U2 – Integración y Pruebas (I\&T / Test Engineer):}
\begin{itemize}
    \item \textbf{Rol:} Prepara el sistema (calibraciones, interfaces), ejecuta pruebas T/A/D/I y documenta evidencias.
    \item \textbf{Características:} Acceso físico al rover; requiere trazabilidad clara entre requisitos y criterios de prueba.
\end{itemize}

\subsection{Stakeholder Needs Register}\label{subsec:stakeholder_needs}

Este registro formaliza las necesidades de cada stakeholder y su derivación hacia requisitos del sistema, conforme a ISO/IEC/IEEE 29148 §9.3 (StRS layer). Cada necesidad tiene un identificador único (SN-xxx), stakeholder fuente, prioridad MoSCoW y los requisitos SyRS/SRS que la satisfacen.

\begin{table}[H]
\centering
\scriptsize
\begin{tabularx}{\textwidth}{| l | l | X | c | p{3.8cm} |}
\hline
\textbf{SN-ID} & \textbf{Stakeholder} & \textbf{Necesidad} & \textbf{MoSCoW} & \textbf{Requisito(s) derivado(s)} \\ \hline

SN-001 & SETEC Lab / TEC &
Validar ELANav en terreno análogo con navegación autónoma real, percepción visual y clasificación de transitabilidad &
Must &
RF-001, RF-002, RF-003, UC-01 \\ \hline

SN-002 & SETEC Lab / TEC &
Demostrar resiliencia del sistema ante fallos de tracción, energía y comunicación sin intervención humana &
Must &
RF-005, SYS-FUN-040a/b, SYS-FUN-041, MODE-005, MODE-006 \\ \hline

SN-003 & Investigador Principal &
Trazabilidad completa desde necesidades operacionales hasta evidencia de V\&V, sin brechas en el digital thread del proyecto &
Must &
§12 Traceability, RTM (Tab.~\ref{tab:rtm}), §13 CR Register \\ \hline

SN-004 & Investigador Principal &
Odometría verificable en campo con error cuantificable respecto a ground truth físico &
Must &
RF-004-R1, RNF-003, GNC-REQ-003, CR-002 \\ \hline

SN-005 & Equipo I\&T &
Sistema desplegable en el sitio análogo en tiempo operacionalmente aceptable para campañas de campo &
Should &
USA-REQ-002, PKG-REQ-001, LCS-REQ-001 \\ \hline

SN-006 & Equipo I\&T &
Telemetría y logs de misión persistentes y accesibles para análisis post-campaña &
Must &
CDH-REQ-002, INF-REQ-001, INF-REQ-002, RNF-007 \\ \hline

SN-007 & Evaluadores / Comité &
Documentación reproducible con evidencia objetiva y trazable para defensa académica del TFG &
Must &
§11 V\&V, §13 FMEA (Tab.~\ref{tab:fmea}), RTM (Tab.~\ref{tab:rtm}) \\ \hline

SN-008 & Evaluadores / Comité &
Integridad de datos de misión verificable mediante mecanismos de validación documentados &
Should &
RF-006, COMM-REQ-001, COMM-REQ-004 \\ \hline

SN-009 & Equipo ELANav &
Plataforma rover capaz de hospedar y validar la biblioteca embebida basada en agentes, exponiendo interfaces de comandos/telemetría y captura de evidencia en sitio análogo &
Must &
UC-01, RF-001, RF-002, RF-003, RF-004-R1, RF-005, RF-006, COMM-REQ-001, GNC-REQ-001, GNC-REQ-002, GNC-REQ-003 \\ \hline

\end{tabularx}
\caption{Registro de necesidades de stakeholders (StRS layer — ISO 29148 §9.3).}
\label{tab:sn_register}
\end{table}

\subsection{Actors and External Systems}\label{subsec:actors_systems}

\textbf{Actores principales (interacción directa):}
\begin{itemize}
    \item \textbf{A1) Rover Planner / Estación Terrena (GCS):} Envía comandos (uplink) y recibe telemetría (downlink) soportando el ciclo diario operacional.
    \item \textbf{A2) Entorno (Environment / Terreno Análogo):} Provee condiciones físicas (pendientes, obstáculos) y contingencias (eventos de deslizamiento/stall de encoder) que disparan comportamientos definidos en \S\ref{sec:requirements}.
\end{itemize}

\textbf{Sistemas externos de soporte:}
\begin{itemize}
    \item \textbf{E1) Infraestructura de comunicaciones:} Soporta el intercambio TM/CMD y el encapsulamiento CSP; su desempeño afecta la latencia.
    \item \textbf{E2) Instrumentación del sitio:} Provee mediciones auxiliares y soporte de registro durante las pruebas del UC‑01.
\end{itemize}
